\documentclass[11pt, letterpaper]{article}

% ── Packages ──────────────────────────────────────────────────────────────────
\usepackage[margin=1.1in, top=1in, bottom=1in]{geometry}
\usepackage{fontspec}                  % requires XeLaTeX or LuaLaTeX
\usepackage{microtype}
\usepackage{setspace}
\usepackage{parskip}
\usepackage{titlesec}
\usepackage{titling}
\usepackage{enumitem}
\usepackage{booktabs}
\usepackage{array}
\usepackage{longtable}
\usepackage{tabularx}
\usepackage{xcolor}
\usepackage{hyperref}
\usepackage{graphicx}
\usepackage{float}
\usepackage{caption}
\usepackage{fancyhdr}
\usepackage{mdframed}
\usepackage{soul}
\usepackage{footnote}
\usepackage{multirow}

% ── Fonts ─────────────────────────────────────────────────────────────────────
\setmainfont{Inter}[
  BoldFont    = Inter Bold,
  ItalicFont  = Inter Italic,
]
% Fallback if Inter not installed: comment above and uncomment below
% \usepackage{helvet}
% \renewcommand{\familydefault}{\sfdefault}

% ── Colors ────────────────────────────────────────────────────────────────────
\definecolor{navy}{HTML}{1a1a2e}
\definecolor{accent}{HTML}{2d6bff}
\definecolor{lightgray}{HTML}{f5f5f7}
\definecolor{midgray}{HTML}{8e8e93}
\definecolor{rulegray}{HTML}{d1d1d6}

% ── Hyperref ──────────────────────────────────────────────────────────────────
\hypersetup{
  colorlinks  = true,
  linkcolor   = accent,
  urlcolor    = accent,
  citecolor   = accent,
  pdftitle    = {BSV AI Fellow -- Written Project},
  pdfauthor   = {Your Name},
}

% ── Section Styling ───────────────────────────────────────────────────────────
\titleformat{\section}
  {\large\bfseries\color{navy}}
  {}
  {0em}
  {}
  [\vspace{-0.3em}\textcolor{rulegray}{\rule{\linewidth}{0.4pt}}\vspace{0.3em}]

\titleformat{\subsection}
  {\normalsize\bfseries\color{navy}}
  {}
  {0em}
  {}

\titleformat{\subsubsection}
  {\normalsize\bfseries\color{midgray}}
  {}
  {0em}
  {}

\titlespacing{\section}{0pt}{1.6em}{0.5em}
\titlespacing{\subsection}{0pt}{1.2em}{0.3em}
\titlespacing{\subsubsection}{0pt}{0.8em}{0.2em}

% ── Page Header / Footer ──────────────────────────────────────────────────────
\pagestyle{fancy}
\fancyhf{}
\renewcommand{\headrulewidth}{0pt}
\fancyhead[L]{\small\textcolor{midgray}{BSV AI Fellow -- Written Project}}
\fancyhead[R]{\small\textcolor{midgray}{Confidential}}
\fancyfoot[C]{\small\textcolor{midgray}{\thepage}}

% ── List Styling ──────────────────────────────────────────────────────────────
\setlist[itemize]{leftmargin=1.4em, itemsep=0.2em, topsep=0.3em}
\setlist[enumerate]{leftmargin=1.4em, itemsep=0.2em, topsep=0.3em}

% ── Callout Box ───────────────────────────────────────────────────────────────
\newmdenv[
  backgroundcolor  = lightgray,
  linecolor        = rulegray,
  linewidth        = 0.5pt,
  roundcorner      = 4pt,
  innerleftmargin  = 14pt,
  innerrightmargin = 14pt,
  innertopmargin   = 10pt,
  innerbottommargin= 10pt,
  skipabove        = 0.8em,
  skipbelow        = 0.8em,
]{callout}

% ── Caption Styling ───────────────────────────────────────────────────────────
\captionsetup{
  font      = {small, color=midgray},
  labelfont = bf,
  skip      = 4pt,
}

% ── Line Spacing ──────────────────────────────────────────────────────────────
\setstretch{1.25}

% ─────────────────────────────────────────────────────────────────────────────
\begin{document}

% ── Title Block ───────────────────────────────────────────────────────────────
\begin{center}
  {\fontsize{22}{26}\selectfont\bfseries\color{navy} BSV AI Fellow}\\[0.4em]
  {\fontsize{13}{16}\selectfont\color{midgray} Written Project Submission}\\[1.2em]
  {\normalsize\textbf{Your Name}} \quad\textcolor{rulegray}{|}\quad
  {\normalsize\textcolor{midgray}{your@email.com}} \quad\textcolor{rulegray}{|}\quad
  {\normalsize\textcolor{midgray}{May 2026}}
\end{center}

\vspace{0.5em}
\textcolor{rulegray}{\rule{\linewidth}{0.6pt}}
\vspace{1.2em}

% ─────────────────────────────────────────────────────────────────────────────
%  PART 1 -- PRIORITIZATION FRAMEWORK
% ─────────────────────────────────────────────────────────────────────────────
\section{Part 1 \quad Prioritization Framework}

\subsection{Evaluation Approach}

For each company I started by understanding what the business actually does and
what pain its customers are experiencing. I watched product demos, reviewed
documentation, and where possible tested the product directly. From there I
mapped the competitive landscape to assess whether the insight was differentiated
or replicable. Finally I looked at the financials, funding history, and founder
background to stress-test whether the team was credible and the business had
real traction behind it. I weighted technical moat and founder-market fit most
heavily, treating traction signals and capital efficiency as confirming or
disconfirming evidence rather than leading indicators.

\subsubsection{Scoring Criteria}

I evaluated each company across four dimensions:

\begin{enumerate}
  \item \textbf{Technical Moat} -- Does the company have a defensible insight,
    proprietary algorithm, or infrastructure advantage that is hard to replicate?
    Preference for companies where the core IP is the product, instead of a wrapper.

  \item \textbf{Founder-Market Fit} -- Does the founding team have direct, earned
    experience with the problem? Operator founders who have experienced the pain
    points are weighted above general builders.

  \item \textbf{TAM \& Timing} -- Is the market large enough to support a
    venture-scale outcome, and is the problem becoming urgent now due to a
    structural shift such as the rise of agentic AI workflows?

  \item \textbf{Traction Signal} -- Is there evidence of real pull from developer
    adoption, revenue growth, or enterprise contracts? Organic traction is
    weighted above funded distribution.
\end{enumerate}

\subsection{Company Rankings}

\begin{longtable}{@{} p{0.04\linewidth} p{0.10\linewidth} p{0.07\linewidth}
                      p{0.28\linewidth} p{0.42\linewidth} @{}}
  \toprule
  \textbf{\#} &
  \textbf{Company} &
  \textbf{Stage} &
  \textbf{Description} &
  \textbf{Reasoning} \\
  \midrule
  \endfirsthead

  \multicolumn{5}{l}{\small\textit{continued from previous page}} \\
  \toprule
  \textbf{\#} & \textbf{Company} & \textbf{Stage} &
  \textbf{Description} & \textbf{Reasoning} \\
  \midrule
  \endhead

  \bottomrule
  \endfoot

  1 & Parasail & Series A &
  AI inference cloud that aggregates GPU supply across 40+ global data centers
  with pay-per-token pricing and no contracts. &
  Up to 30x cheaper than legacy cloud with 500B tokens processed per day and
  30\% MoM revenue growth less than a year post-launch. Faces competition from
  Fireworks and Baseten. CEO Mike Henry built Groq Cloud. Co-founder Tim Harris
  raised \$250M at Swift Navigation. \\[0.5em]

  2 & Tigris Data & Series A &
  Globally distributed object storage built for AI workloads that eliminates
  cloud egress fees by moving data to where GPUs are. &
  Egress fees accounted for the majority of fal.ai cloud spend before switching
  to Tigris. Zero-copy Forks and data gravity create stickiness in a
  fast-growing TAM with 8x YoY growth. Founded by Ovais Tariq, who built and
  scaled Uber's global storage platform. \\[0.5em]

  3 & Entire & Seed &
  Developer platform that captures agent session context and ties it into Git,
  creating a persistent reasoning layer across agentic coding workflows. &
  Tools like Claude Code are making multi-agent workflows the default, yet agent
  sessions remain temporary with no shared context. High displacement risk from
  GitHub if they move into this space. \$60M seed backed by Garry Tan, Olivier
  Pomel (Datadog CEO), and more, signaling strong operator conviction. \\[0.5em]

  4 & BAML & Seed &
  Domain-specific language for building type-safe LLM applications that handles
  structured output parsing, retries, and testing across any model or language. &
  Addresses enterprise LLM reliability via cost-advantaged Schema-Aligned
  Parsing. Outperforms GPT-4o native outputs using GPT-3.5. High schema
  stickiness and cross-modal portability provide a strong moat against native
  rivals like Anthropic Tool Use. \\[0.5em]

  5 & Mem0 & Seed &
  Persistent memory layer that gives AI agents and applications long-term recall
  across sessions, users, and agents. &
  Making persistent memory a prerequisite for any long-running agent product.
  Large and growing TAM as agent adoption accelerates. Faces the risk of being
  commoditized by model providers and competition from well-funded players
  like Zep. \\[0.5em]

  6 & Rasa & Series C &
  Open-source enterprise conversational AI platform that combines LLM fluency
  with controllable business logic, deployed by Fortune 500 companies at scale. &
  Proven enterprise traction with 50M+ downloads and PayPal deploying
  internally, but open source to enterprise conversion at scale remains
  questionable. CALM architecture is a real moat, however Sierra, Cognigy, and
  native agent frameworks are closing the gap fast at a late stage with limited
  new diligence signal to uncover. \\[0.5em]

  7 & Atuin & Pre-Seed &
  Replaces shell history with an encrypted database synced across machines,
  expanding into executable runbooks via Atuin Desktop. &
  Data gravity and a loyal open source community trusted by thousands of
  developers worldwide. The pivot to executable runbooks is the more interesting
  enterprise bet. Monetization path is unclear and not an obvious venture-scale
  outcome without the runbook vision materializing. \\[0.5em]

  8 & Primitive & Seed &
  Email infrastructure for developers and AI agents, offering managed webhooks
  and self-hosted options with SDKs for Node.js, Python, and Go. &
  Abstracts complex email infrastructure into a developer-friendly API with a
  time-to-market advantage in the AI agent space. High infrastructure costs at
  scale with unproven IP reputation against Resend and SendGrid make enterprise
  adoption difficult. \\

\end{longtable}

\subsection{Supporting Data}

\subsubsection{Funding Overview}

{\small\textcolor{midgray}{Source: PitchBook, TechCrunch, Tracxn, company press releases.}}

\vspace{0.4em}

\begin{table}[H]
\centering
\begin{tabularx}{\linewidth}{@{} l r l X @{}}
  \toprule
  \textbf{Company} & \textbf{Total Raised} & \textbf{Last Round} & \textbf{Key Signal} \\
  \midrule
  Entire      & \$60.0M               & Seed, Feb 2026     & Largest developer tools seed ever, \$300M valuation \\
  Rasa        & \$83.8M               & Series C, Feb 2024 & \$16.9M ARR in 2023, 50M+ downloads across 9 rounds \\
  Parasail    & \$42.0M               & Series A, Apr 2026 & 500B tokens/day, 30\% MoM revenue growth \\
  Tigris Data & \$39.2M               & Series A, Oct 2025 & 8x YoY growth, 4,000+ customers \\
  Mem0        & \$24.5M               & Series A, Oct 2025 & 41K GitHub stars, 186M API calls in Q3 2025 \\
  BAML        & \textasciitilde\$130K & Seed, undisclosed  & Most capital-efficient company on this list \\
  Atuin       & Undisclosed           & Pre-Seed           & No disclosed amount \\
  Primitive   & Undisclosed           & Seed, Apr 2026     & Launched Apr 2026, no disclosed amount \\
  \bottomrule
\end{tabularx}
\end{table}

\begin{callout}
  \textbf{Insight.} BAML has raised approximately \$130K yet commands the
  strongest organic developer community on this list. Funding size does not
  correlate with diligence priority. It is one signal among four.
\end{callout}

\subsubsection{LinkedIn Headcount}

{\small\textcolor{midgray}{Source: LinkedIn Premium, May 2026.}}

\vspace{0.4em}

\begin{table}[H]
\centering
\begin{tabularx}{\linewidth}{@{} l r r r r r X @{}}
  \toprule
  \textbf{Company} &
  \textbf{Employees} &
  \textbf{6m} &
  \textbf{1y} &
  \textbf{2y} &
  \textbf{Tenure} &
  \textbf{Signal} \\
  \midrule
  Entire    & 28  & +211\% & ---   & ---    & 0.3yr &
    211\% surge in 6m post \$60M seed. Capital being deployed fast. \\[0.3em]
  Mem0      & 21  & +62\%  & +110\% & +600\% & 0.4yr &
    600\% over 2y is the most aggressive hiring trajectory on this list \\[0.3em]
  Parasail  & 31  & +35\%  & +82\%  & +158\% & 0.9yr &
    Consistent acceleration across all time horizons confirms real revenue pull \\[0.3em]
  BAML      & 11  & +57\%  & +22\%  & +120\% & 0.3yr &
    11 people with outsized community traction. Capital efficiency thesis confirmed. \\[0.3em]
  Rasa      & 194 & +22\%  & +34\%  & +59\%  & 1.4yr &
    Largest and most stable team. Slowest growth rate consistent with Series C maturity. \\[0.3em]
  Tigris    & 23  & +10\%  & +10\%  & +77\%  & 1.8yr &
    Flat near-term growth despite fresh Series A. Warrants a diligence question. \\[0.3em]
  Primitive & 3   & +200\% & ---    & ---    & 0.2yr &
    200\% in 4m but from a base of 1. Too early to read meaningfully. \\[0.3em]
  Atuin     & 2   & 0\%    & 0\%    & +100\% & 1.6yr &
    Zero hiring in 12 months. Solo founder mode limits the enterprise runbook thesis. \\
  \bottomrule
\end{tabularx}
\end{table}

\begin{callout}
  \textbf{Insight.} Tigris shows only +10\% headcount growth over 12 months
  despite claiming 8x YoY revenue growth and closing a fresh Series A. The
  disconnect warrants a direct diligence question. Either growth is concentrated
  in a small number of large contracts, or the revenue figure is being measured
  differently than it appears.
\end{callout}

% ─────────────────────────────────────────────────────────────────────────────
%  PART 2 -- TRACK 2: AGENTIC GTM -- VIZCOM
% ─────────────────────────────────────────────────────────────────────────────
\newpage
\section{Part 2 \quad Agentic GTM -- Vizcom}

\subsection{Understanding the Company}

\subsubsection{What Vizcom Does}

Vizcom turns rough sketches into photorealistic 3D renders in seconds, built
for physical product designers rather than digital ones. A designer draws in
Procreate or on paper, uploads it, and gets a material-accurate product
visualization in real time with no rendering artist, no outsourcing, and no
KeyShot export queue. The CEO is an ex-Honda, ex-Nvidia industrial designer.
The product writes like someone who has actually used KeyShot and found it
exhausting.

\subsubsection{The Growth Problem}

Design agencies win or lose pitches based on how many concept directions they
can show. Traditional rendering creates a hard ceiling. Teams sketch 50 ideas,
render 4, and the client picks from 4. Vizcom removes that ceiling. The growth
problem is awareness. Most boutique agencies have never heard of it, and the
buyer who controls budget is not the same person who
would stumble across it organically.

\subsection{GTM Agent Design}

\subsubsection{What the Agent Does}

The agent monitors the web for signals that a boutique industrial design agency
is actively pitching new business. When it detects a qualifying signal, it
researches the agency, scores the lead against a weighted rubric, and generates
personalized outreach across a coordinated channel sequence. It tracks
responses and retrains its scoring weights weekly based on reply rates.

\subsubsection{Inputs and Outputs}

\begin{table}[H]
\centering
\begin{tabularx}{\linewidth}{@{} l X @{}}
  \toprule
  & \textbf{Description} \\
  \midrule
  \textbf{Inputs} &
  LinkedIn job postings, Behance and ArtStation portfolio projects,
  Core77 and Dezeen trade press, a seed list of 500+ boutique ID consultancies,
  and LinkedIn employee profiles for tool stack signals (KeyShot, Rhino, Alias). \\[0.5em]
  \textbf{Outputs} &
  Lead scorecard (JSON), cold email with two subject line variants, Behance
  portfolio comment (pre-email, no pitch), LinkedIn connection
  note and DM sequence, physical mailer for 80+ score leads
  via Lob API, CRM log entry, weekly performance digest. \\
  \bottomrule
\end{tabularx}
\end{table}

\subsubsection{Lead Scoring}

Leads are scored 0 to 100. Scores above 65 proceed to outreach. Scores 40
to 64 enter a watch list.

\begin{table}[H]
\centering
\begin{tabularx}{\linewidth}{@{} l r X @{}}
  \toprule
  \textbf{Dimension} & \textbf{Weight} & \textbf{Logic} \\
  \midrule
  Vertical fit    & 25 pts & Footwear, electronics, furniture = 25. Automotive = 15. Medical = 5. \\
  Company size    & 20 pts & 10 to 25 employees = 20. Outside range = 0. \\
  Tool signal     & 20 pts & KeyShot or Rhino in employee profiles = 20. No signal = 5. \\
  Signal recency  & 20 pts & Under 7 days = 20. 7 to 14 days = 12. Over 14 days = 5. \\
  Signal strength & 15 pts & New client win = 15. Case study = 12. Job post = 8. Portfolio update = 5. \\
  \bottomrule
\end{tabularx}
\end{table}

\subsubsection{Channel Sequence}

Channels run in a coordinated sequence, not in parallel. Each touch warms the next.

\begin{table}[H]
\centering
\begin{tabularx}{\linewidth}{@{} l l X @{}}
  \toprule
  \textbf{Day} & \textbf{Channel} & \textbf{Action} \\
  \midrule
  0       & Detection          & Agency posts case study, new project, or junior designer listing \\
  1       & Behance/ArtStation & Opinionated comment on their new project. No pitch, no link. \\
  3       & Cold email         & Sent to Design Director. References specific project. Closes with a question. \\
  7       & LinkedIn           & Connection request with short note referencing the same project \\
  14      & LinkedIn DM        & Follow-up sharing something useful in their vertical. No hard ask. \\
  25--28  & Physical mailer    & 80+ score leads only. Printed before/after render via Lob API. \\
  30      & Final email        & Short. Acknowledges they may not be the right fit. Leaves the door open. \\
  \bottomrule
\end{tabularx}
\end{table}

\subsection{Technical Stack}

\begin{table}[H]
\centering
\begin{tabularx}{\linewidth}{@{} l X @{}}
  \toprule
  \textbf{Component} & \textbf{Tool} \\
  \midrule
  Orchestration       & Python with APScheduler (daily detection, weekly learning loop) \\
  LinkedIn signals    & LinkedIn Jobs API and Apify scraper \\
  Portfolio signals   & Apify actor for Behance. ArtStation disabled (no reliable actor). \\
  Trade press         & RSS feeds for Core77, Dezeen, and Google News \\
  LLM classification  & \texttt{claude-haiku-4-5} (approx. 200 calls/day) \\
  LLM scoring/brief   & \texttt{claude-haiku-4-5} (approx. 30 calls/day) \\
  LLM email/DM copy   & \texttt{claude-sonnet-4-20250514} (approx. 30 calls/day) \\
  Database            & SQLite (dev) to PostgreSQL (prod) \\
  Email sending       & SendGrid on a warmed subdomain \\
  LinkedIn outreach   & Phantombuster within ToS rate limits \\
  Physical mail       & Lob API (not yet implemented) \\
  CRM                 & HubSpot API or CSV export (not yet implemented) \\
  Source code         & \href{https://github.com/bhargavap21/outboundvizcom}{github.com/bhargavap21/outboundvizcom} \\
  \bottomrule
\end{tabularx}
\end{table}

\subsection{Example Output}

\subsubsection{Cold Email -- Whipsaw (Consumer Electronics)}

\begin{callout}
  \textbf{Subject A:} How many renders were in your last pitch deck?\\
  \textbf{Subject B:} You sketched 30. You showed 4.\\[0.8em]
  Dan,\\[0.5em]
  The temple-to-cup junction on your headphone project is the kind of decision
  that only gets made by someone who has actually held a bad pair. It is
  considered in a way most concept work is not.\\[0.5em]
  Quick question: how many render variations did you show the client before
  landing on that direction?\\[0.5em]
  Every studio I talk to has the same problem. You sketch 30 directions, you
  render 4, and the client picks from 4. The constraint is never ideas.\\[0.5em]
  New Balance's footwear team cut their colorway review cycle in half.
  Breville uses it before committing to a single physical sample.\\[0.5em]
  Curious whether your render cycle is actually the bottleneck, or whether
  it sits somewhere else in your process.\\[0.5em]
  [Name]
\end{callout}

\subsection{Success Metrics}

\begin{table}[H]
\centering
\begin{tabularx}{\linewidth}{@{} X r r @{}}
  \toprule
  \textbf{Metric} & \textbf{Month 1} & \textbf{Month 3} \\
  \midrule
  Qualified leads per week          & 15--20   & 40--60 \\
  Email open rate                   & >35\%    & >45\% \\
  Email reply rate                  & >5\%     & >8\% \\
  LinkedIn connection acceptance    & >30\%    & >40\% \\
  Demos booked per month            & 3--5     & 10--15 \\
  Cost per qualified lead           & <\$2.00  & <\$1.00 \\
  \bottomrule
\end{tabularx}
\end{table}

\subsection{Evaluation}

\subsubsection{What Was Built and What Worked}

The core signal-to-content pipeline is fully operational end-to-end. Signal
detection pulls live data from three sources. LinkedIn posts via Apify, Google
News RSS (Core77, Dezeen, Designboom), and Behance portfolio updates via the
\texttt{easyapi/behance-people-scraper} Apify actor. In a test run against the
seed agency list, the pipeline surfaced real signals for Morrama, Priority
Designs, Whipsaw, Bould Design, and Layer within a single daily pass, all
verified against live URLs.

The scoring rubric runs as designed. Every signal is classified by Claude Haiku
(vertical, signal type, size hint), enriched with team size from the seed
database, and scored against the five-dimension rubric. Morrama and Priority
Designs both scored 82 on case studies in relevant verticals. Whipsaw scored 65
on a client win in consumer electronics. Haiku classification correctly
identified verticals and signal types across every tested agency with no
manual review needed.

Content generation produces all five outreach pieces per lead. Behance comment,
cold email with two subject line variants, LinkedIn connection note, Day-14 DM,
and Day-30 final email. The voice constraints are enforced at the prompt level.
The banned word list covers ``AI,'' ``seamless,'' ``game-changing,'' and
eleven other terms. Every email in test output opened with a specific
observation about the agency's work and closed with a question rather than a
meeting request.

The distribution layer is wired together. The sequence dispatcher advances
leads through timed steps (Day 3, 7, 14, 30) using a database-backed idempotency
check so no step fires twice. The email sender enforces a 40-per-day hard cap
and an 8am--11am send window. The LinkedIn sender enforces the 20-per-day
connection limit. The inbound response handler uses Claude Haiku to classify
replies into five categories (positive, objection, unsubscribe, bounce,
no\_response), drafts objection replies for human review, suppresses
unsubscribes immediately, and posts warm lead alerts to Slack.

APScheduler orchestrates the full loop. A daily job fires signal detection,
scoring, research, content generation, and sequence dispatch. A Sunday job fires
the weekly learning loop, which computes open and reply rates by variant and
produces a markdown digest via Claude Sonnet.

\subsubsection{What Broke or Fell Short}

\textbf{Tool stack detection is effectively broken.} The scoring rubric assigns
20 points for a KeyShot or Rhino signal, but the enricher only scans the text
of incoming signal posts for tool keywords. Agency news posts and case studies
do not mention their renderer. Every lead in the real test run showed no tool
stack detected, meaning the tool dimension contributed 5 points (the minimum)
across 100\% of scored leads. The intended mechanism, scanning LinkedIn
employee profiles for tool keywords, was deferred to the full research step,
which only runs after a lead already passes the threshold. The dimension is
load-bearing in the scoring model but currently has no data to feed it.

\textbf{ArtStation is disabled.} No reliable Apify actor exists for ArtStation
as of May 2026. The \texttt{fetch\_artstation\_projects} function returns an
empty list immediately. Behance works, but the planned two-platform portfolio
signal layer is currently one platform.

\textbf{Behance comments are not automated.} The content generator produces the
comment copy for each lead, but there is no mechanism to post it. The sequence
dispatcher marks Day-1 Behance comments as ``manual or future automation.'' The
content exists in the database. The posting step does not.

\textbf{Physical mailers are not built.} The spec calls for Lob API dispatches
for 80+ score leads on Day 25--28. The sequence dispatcher's step schedule
contains no Lob step and there is no Lob client anywhere in the codebase. This
is the largest unbuilt component.

\textbf{The inbound response handler has no entry point.} The handler correctly
classifies and routes replies, but there is no HTTP webhook endpoint for
SendGrid to call when a reply arrives. In production it needs a lightweight
FastAPI server wired to SendGrid's Inbound Parse webhook.

\textbf{The learning loop does not retrain weights.} The weekly analyzer
computes open and reply rates by variant and generates a human-readable markdown
digest, but it does not update the scorer's dimension weights. The weights are
hardcoded in \texttt{scorer.py}. The spec's claim that the agent ``retrains its
scoring weights weekly'' is not implemented.

\textbf{LinkedIn acceptance is inferred, not verified.} The
\texttt{is\_connected()} check confirms only that a connection request was
\emph{dispatched}, not that it was accepted. Day-14 DMs gate on this check, so
the agent cannot confirm actual connection status before sending.

\textbf{The live pipeline runs against an empty seed list.} The \texttt{SEED\_AGENCIES}
list in \texttt{main.py} is hardcoded as empty with a \texttt{TODO} comment to
load from the database. The daily job would detect zero signals on a live run.
The seed list exists in the \texttt{AgencyList} table but the main loop does
not load from it.

\subsubsection{What I Would Do With More Time}

\begin{itemize}
  \item \textbf{Fix the seed list load.} One line. Replace the empty
    \texttt{SEED\_AGENCIES} list in \texttt{main.py} with a DB query against
    \texttt{AgencyList}. This is the highest-leverage unblocked fix.

  \item \textbf{Move tool stack detection to pre-scoring.} The enricher should
    scrape LinkedIn employee profiles for each agency before scoring, not after.
    The tool signal dimension is the most predictive in the rubric and currently
    contributes nothing. Running the Apify LinkedIn Profiles Scraper on the top
    3 employees per agency would restore the signal before the threshold
    decision is made.

  \item \textbf{Build the Lob integration.} The physical mailer is the most
    differentiated component in the channel sequence and the only piece with no
    equivalent in a standard outbound stack. The 80+ score filter keeps volume
    low enough (2--5 per week) that unit economics are manageable from day one.

  \item \textbf{Expose the response handler as a webhook.} A single FastAPI
    endpoint wired to SendGrid's Inbound Parse makes the response loop live.
    Without it, the objection drafts and Slack alerts the handler generates are
    never triggered.

  \item \textbf{Implement actual weight regression in the learning loop.} The
    analyzer already collects the right data (reply rates by scoring dimension).
    The missing piece is a regression step that updates the weights in
    \texttt{scorer.py} based on six weeks of outcome data. Company size is the
    weakest prior and needs the most empirical validation.

  \item \textbf{Add a vision model pass on Behance renders.} Agencies whose
    public renders look rough are higher-priority targets than agencies whose
    renders already look polished. A multimodal pass on Behance project images
    for each 65+ score lead would add a sixth scoring dimension unavailable to
    any competitor using the same text-only pipeline.
\end{itemize}

% ─────────────────────────────────────────────────────────────────────────────
%  FOOTER NOTE
% ─────────────────────────────────────────────────────────────────────────────
\vfill
\textcolor{rulegray}{\rule{\linewidth}{0.4pt}}
\begin{center}
  \small\textcolor{midgray}{
    LLM conversations are attached separately as supplementary material.
    All funding data sourced from PitchBook, TechCrunch, Tracxn, and company
    press releases as of May 2026.
  }
\end{center}

\end{document}
